\documentclass[12pt,a4paper]{article}

\usepackage[utf8]{inputenc}
\usepackage[T1]{fontenc}
\usepackage{lmodern}
\usepackage[margin=1in]{geometry}
\usepackage{setspace}
\usepackage[hidelinks]{hyperref}
\usepackage{amsmath,amssymb}
\usepackage{enumitem}
\usepackage{booktabs}
\usepackage{titlesec}
\usepackage[stretch=15]{microtype}
\usepackage{parskip}
\usepackage{longtable}
\usepackage{array}

% Section formatting
\titleformat{\section}{\large\bfseries}{}{0em}{}
\titleformat{\subsection}{\normalsize\bfseries}{}{0em}{}
\titleformat{\subsubsection}{\normalsize\itshape}{}{0em}{}

\title{Technical Review Memo\\[6pt]
\large ``Principled Hyperprior Selection for Bayesian Hierarchical Models\\
of Risk Preferences''\\[6pt]
\normalsize Target Journal: \textit{Journal of Business \& Economic Statistics}}
\author{}
\date{February 2026 (Updated)}

\begin{document}
\maketitle
\thispagestyle{empty}

\vspace{1em}
\hrule
\vspace{1em}

\section{Current State of the Paper}

The paper has undergone significant development since the initial review.  The
major structural gaps identified previously---an empty Appendix~A and a sparse
Appendix~C---have been resolved.  Prose has been tightened throughout, with
consistent first-person singular usage, improved figure captions, and more
precise language in the Results and Discussion sections.

\textbf{Resolved since last review:}
\begin{itemize}[leftmargin=1.5em]
\item \textbf{Appendix~A} (Full Covariance Derivation) is now complete.  It
provides the full Normal-Inverse-Wishart extension: setup, NIW hierarchical
structure, mean and covariance constraints via matrix law of total variance,
variance partition with $\rho$, derived hyperparameters
($\boldsymbol{\mu}_0$, $\kappa_0$, $\boldsymbol{\Psi}_0$, $\nu_0$),
verification that the unconditional covariance recovers $\boldsymbol{\Sigma}^*$,
the relationship to the independent case (showing IW reduces to independent
IG), and a note on Cholesky parameterization for sampling.

\item \textbf{Appendix~C} (MCMC Implementation Details) has been expanded from
3 short paragraphs to 9 substantive subsections covering: the sampling
algorithm (115 blocks, 452 parameters, coupled A/B chains), non-centered
parameterization with funnel motivation, the coupled transform sampler with the
full log-acceptance formula (Jacobian, epsilon prior ratio, likelihood
cancellation), proposal adaptation via cross-chain covariance, run
configuration (1000 chains, 6 runs, thinning, burn-in), convergence diagnostics
($\hat{R}$), and JAX 64-bit precision.

\item \textbf{Prose consistency:} ``we'' $\to$ ``I'' throughout, ``uninformative''
$\to$ ``diffuse'' where appropriate, section titles improved (e.g.,
``Hyper-parameter Comparison'' $\to$ ``Hyperparameter Posteriors''), more precise
descriptions of figures and results.
\end{itemize}

\textbf{The paper now has a complete technical core.}  The remaining gaps are
about positioning, robustness evidence, and the additional empirical content
that JBES referees will expect.

\vspace{0.5em}
\hrule
\vspace{0.5em}

\section{1.\ Issues for JBES Submission}

The following issues are organized by priority for a JBES submission.  Issues
that were resolved are not repeated; only remaining gaps are listed.

\subsection{1.1\quad Literature Positioning [Critical]}

This remains the most critical gap.  The paper cites only Gao, Harrison \&
Tchernis (2022) as the main comparator and does not engage with the substantial
existing literature on hyperprior selection for hierarchical models.  A JBES
referee will expect engagement with several strands of work.

\subsubsection{Variance components and hyperprior pathologies.}

The canonical reference is \textbf{Gelman (2006), ``Prior distributions for
variance parameters in hierarchical models''} (\textit{Bayesian Analysis},
1(3), 515--534).  Gelman argues against Inverse-Gamma priors for variance
components and proposes half-Cauchy alternatives, demonstrating that the common
``non-informative'' choice $\text{IG}(\varepsilon, \varepsilon)$ concentrates
mass near zero.  Since the paper's framework is built on Normal--Inverse-Gamma
conjugacy, it must engage with this critique and explain why the proposed
variance decomposition avoids these pathologies.  The key defense---that the
proposed method \emph{never} uses $\text{IG}(\varepsilon, \varepsilon)$ but
instead derives $\alpha_0$ and $\beta_0$ from substantive beliefs---should be
made explicit.

Follow-ups: \textbf{Polson \& Scott (2012)} on half-Cauchy priors, and
\textbf{Huang \& Wand (2013)} on half-$t$ priors for covariance matrices.
Both offer alternative solutions to the same problem.

\textbf{Berger, Strawderman \& Tang (2005)} establishes when hierarchical
priors yield proper posteriors---directly motivating a principled approach.

\subsubsection{Principled prior construction frameworks.}

\textbf{Simpson et al.\ (2017), ``Penalising model component complexity''}
proposes PC priors that penalize deviations from a base model.  The variance
decomposition has conceptual kinship with this framework---both derive priors
from interpretable principles.  A direct comparison would clarify the
contribution.

\textbf{Natarajan \& Kass (2000)} derives reference priors for variance
components in GLMMs.  How does the paper's approach compare to formal reference
priors?

\textbf{Daniels \& Kass (1999)} develops hierarchical priors for covariance
matrices---now directly relevant to the completed Appendix~A.

\subsubsection{Prior elicitation and predictive checking.}

\textbf{Gabry et al.\ (2019)} establishes prior predictive checking as a core
Bayesian workflow component.  The paper effectively uses prior predictive
distributions (the CE analysis) but does not connect this to the broader
literature.

\textbf{Mikkola et al.\ (2023)} reviews modern prior elicitation methods.  The
paper's approach---eliciting individual-level beliefs and transforming to
population-level hyperpriors---is a form of structured elicitation that should
be positioned within this framework.

\subsubsection{BHMs in experimental economics.}

Beyond Gao et al.\ (2022): \textbf{Nilsson, Rieskamp \& Wagenmakers (2011)} on
hierarchical CPT estimation; \textbf{Conte, Hey \& Moffatt (2011)} on mixture
models for preference heterogeneity; \textbf{Hey \& Orme (1994)} on the
experimental paradigm.

\subsubsection{Model comparison alternatives.}

The Discussion critiques Bayes factors.  This should engage with
\textbf{Vehtari, Gelman \& Gabry (2017)} on LOO-CV as a practical alternative.

\medskip

Without this literature engagement, a JBES editor cannot assess novelty.  The
contribution must be framed as the specific reparameterization and its
interpretability---argued clearly against alternatives.

\subsection{1.2\quad Sensitivity Analysis [Critical]}

The paper claims $\rho$ and $\alpha_0$ are interpretable and that researchers
can meaningfully choose them, but shows only one configuration ($\rho = 0.2$,
$\alpha_0 = 20$).  A JBES methods paper must include:

\begin{itemize}[leftmargin=1.5em]
\item A grid of $(\rho, \alpha_0)$ values---e.g., $\rho \in \{0.05, 0.1, 0.2,
0.5, 0.8\}$ and $\alpha_0 \in \{3, 10, 20, 50, 100\}$---showing how posterior
summaries change.

\item At minimum, a table showing hyperparameter posterior summaries and
subject-level CE statistics for 3--4 different $(\rho, \alpha_0)$ pairs.

\item This addresses the obvious referee question: ``How much do your
conclusions depend on this choice?''  If the answer is ``very little,'' that is
a powerful robustness result.
\end{itemize}

\textbf{Note:} Given the PythonTeX pipeline and the already-working sampling
infrastructure, this is computationally expensive but structurally
straightforward.  The sensitivity analysis does not require new code
architecture---just additional sampling runs with different prior
configurations.

\subsection{1.3\quad Single Dataset [Important]}

For a JBES methods paper, relying on a single empirical application with 111
subjects is thin.  Two potential additions:

\subsubsection{Simulation study.}
Generate data from a known DGP with known population parameters.  Apply the
method with correct and incorrect direct priors.  Show recovery of true
parameters (bias, RMSE, coverage), comparison against diffuse priors, and
degradation under misspecification.

\subsubsection{Second empirical application.}
The Andersen et al.\ (2008) dataset is already implemented in the code
pipeline.  Running the same analysis on this dataset (501 subjects, binary
lotteries, DKK currency) and showing consistent methodology behavior would
strengthen the empirical contribution.

\medskip

Either one of these would substantially strengthen the paper for JBES.  Both
together would make it very competitive.  Of the two, the simulation study is
likely more valued by JBES referees for a methods paper.

\subsection{1.4\quad The Diffuse Prior Comparison [Important]}

The ``diffuse'' prior compared against---$N(0, 100^2)$ for means and
$\text{IG}(0.001, 0.001)$ for variances---is known to be problematic (Gelman,
2006).  Comparing a principled approach against a \emph{known-bad} prior does
not fully demonstrate the method's value.

A fairer comparison would add at least one ``reasonable'' alternative:
\begin{itemize}[leftmargin=1.5em]
\item A half-Cauchy$(0, 2.5)$ prior on $\sigma$ (Gelman's recommendation)
\item A weakly informative but proper prior, e.g., $N(0, 5^2)$ on means,
$\text{IG}(2, 1)$ on variances
\end{itemize}

\textbf{Counterargument:} The diffuse prior from Gao et al.\ (2022) is the one
actually used in the published literature being replicated.  One could argue
that comparing against the prior that was \emph{actually used} is the relevant
comparison, even if that prior is now understood to be suboptimal.  If this
framing is adopted, it should be stated explicitly, and the Gelman (2006)
critique acknowledged.

\subsection{1.5\quad Posterior Summary Table [Important]}

The paper presents results entirely through figures.  A JBES submission needs:
\begin{itemize}[leftmargin=1.5em]
\item A table of posterior means, standard deviations, and 95\% credible
intervals for all 8 hyperparameters under both specifications.

\item Subject-level summary statistics (distribution of posterior mean $r$
across subjects, fraction classified as risk-averse, etc.).

\item Effective sample sizes and $\hat{R}$ statistics in a table, not just a
footnote.
\end{itemize}

\vspace{0.5em}
\hrule
\vspace{0.5em}

\section{2.\ Moderate Issues}

\subsection{2.1\quad Direct Prior Specification Guidance}

The paper still hand-waves past the most important step: ``These prior choices
represent my a-priori belief about what a reasonable subject may look like.''
The paper should:
\begin{itemize}[leftmargin=1.5em]
\item Discuss \emph{how} a researcher should specify the direct prior (published
estimates, meta-analyses, expert elicitation).
\item Provide guidance for common risk preference parameters---a brief survey of
published CRRA $r$ estimates would be valuable.
\item Discuss robustness to moderate misspecification (ties to the simulation
study recommendation).
\end{itemize}

\subsection{2.2\quad Discussion Section Structure}

The Discussion has two points: (1)~method defense and (2)~Bayes factor critique.
The Bayes factor discussion takes up roughly half the section.  For JBES:
\begin{itemize}[leftmargin=1.5em]
\item Consider condensing the Bayes factor critique to one paragraph.
\item Use the freed space for limitations, extensions (e.g., multivariate
case---now supported by Appendix~A), and connections to the broader literature.
\item A proper ``Conclusion'' section (separate from Discussion) would improve
structure.  The paper currently ends abruptly.
\item Mention LOO-CV (Vehtari et al., 2017) as a practical model comparison
alternative when discussing Bayes factor limitations.
\end{itemize}

\subsection{2.3\quad Identification and Data Informativeness}

The paper shows (via the most-affected subject) that prior choice matters when
individual data are uninformative, but provides no formal analysis.  A brief
discussion of:
\begin{itemize}[leftmargin=1.5em]
\item The information content of 80 binary choices for 4 parameters
\item Whether some parameters are more weakly identified ($\phi$ and $\eta$
likely compete through PWF curvature)
\item How partial identification interacts with hierarchical structure
\end{itemize}
would strengthen the results interpretation.

\subsection{2.4\quad Normal--Inverse-Gamma Assumption}

The framework assumes Normal population distributions with IG variance
hyperpriors.  The paper should briefly discuss:
\begin{itemize}[leftmargin=1.5em]
\item Why this is appropriate for transformed RDU parameters (log-transforms of
$\phi$, $\eta$, $\lambda$)
\item The Skew-Normal appendix (B) demonstrates extensibility, but includes no
numerical verification---adding Monte Carlo results would strengthen it
\end{itemize}

\vspace{0.5em}
\hrule
\vspace{0.5em}

\section{3.\ Minor Issues}

\subsection{3.1\quad Writing}

Most of the prose issues from the initial review have been addressed (``we''
$\to$ ``I'', comma splices fixed).  Remaining:
\begin{itemize}[leftmargin=1.5em]
\item ``reign in'' should be ``rein in'' (if still present).
\item Verify ``of during the COVID-19 pandemic'' $\to$ ``during the COVID-19
pandemic'' (extraneous ``of'').
\item Check all \texttt{\textbackslash py\{\}} calls produce sensible output
(not ``0'' or ``N/A'') when the paper compiles with pythontex.
\end{itemize}

\subsection{3.2\quad Notation}

\begin{itemize}[leftmargin=1.5em]
\item The $\sigma^{2*}$ superscript notation becomes confusing with parameter
superscripts in the derived hyperprior equations.  Consider $\sigma^2_d$ or
tilde notation.
\item $\Phi$ is used as the model function in eq.~(1) and as the normal CDF in
the Skew-Normal appendix.  Avoid this collision.
\end{itemize}

\subsection{3.3\quad Figures}

\begin{itemize}[leftmargin=1.5em]
\item Figure captions have been improved (``Hyperparameter posteriors'' rather
than ``Hyper-parameter comparison'') but are still terse.  JBES prefers
self-contained captions (e.g., ``Posterior densities of the 8 hyperparameters
under informative (blue) and diffuse (orange) priors, with the hyperprior
density (grey) shown for reference.  The informative hyperprior is derived via
the variance decomposition in Section~2.'').
\item Ensure figures are legible in grayscale.
\end{itemize}

\subsection{3.4\quad Bibliography}

\begin{itemize}[leftmargin=1.5em]
\item The bibliography style (\texttt{IEEEtran} + \texttt{authoryear-comp}) is
unusual for JBES\@.  Check JBES style requirements---JBES uses Chicago-style
author-year citations.
\item The paper is still light on citations (${\sim}$11 references).  The
literature review expansion (\S1.1) would naturally bring this to 25--30.
\end{itemize}

\subsection{3.5\quad Reproducibility}

\begin{itemize}[leftmargin=1.5em]
\item Add a code/data availability statement.  JBES now requires this.
\item The \texttt{bamcmc} sampler should be cited or described, with a plan for
public availability.
\end{itemize}

\subsection{3.6\quad CE Equation}

Equation~(14) [now the CE equation in the Results section] should be checked for
consistency.  The $(1-r)$ denominators on LHS and RHS should cancel cleanly.
As typeset, the first line retains $1/(1-r)$ on both sides, and the second line
applies the exponent $1/(1-r)$.  Verify the algebra is presented consistently.

\vspace{0.5em}
\hrule
\vspace{0.5em}

\section{4.\ Prioritized Action Items for JBES}

Given that Appendix~A and Appendix~C are now complete, the revised priority
order for a JBES submission is:

\begin{enumerate}[leftmargin=1.5em]
\item \textbf{Literature review and positioning} (\S1.1)---the most likely
reason for a desk reject at JBES.  Without it, the editor cannot assess
novelty.  This is the single highest-priority item.

\item \textbf{Sensitivity analysis} (\S1.2)---the single most impactful
empirical addition.  A table or figure showing posteriors under different
$(\rho, \alpha_0)$ values.  JBES referees will ask for this.

\item \textbf{Posterior summary table} (\S1.5)---basic requirement.  Numbers in
a table, not just curves in figures.

\item \textbf{Expand the diffuse prior comparison} (\S1.4)---either add a
``reasonable'' comparator or explicitly frame the Gao et al.\ prior as the
published benchmark.

\item \textbf{Simulation study or second application} (\S1.3)---either would
elevate the paper.  A simulation study is preferred for a methods-oriented
journal like JBES.

\item \textbf{Discussion restructuring and conclusion} (\S2.2)---condense Bayes
factor critique, add limitations and extensions, add a conclusion section.

\item \textbf{Direct prior specification guidance} (\S2.1)---a paragraph or two
on how researchers should think about this step.

\item \textbf{Figure captions, bibliography style, and remaining minor fixes}
(\S\S3.1--3.6)---low effort, high signal of paper maturity.
\end{enumerate}

\textbf{Minimum for submission:} Items 1--3 are likely necessary to avoid a desk
reject.  Items 4--5 would be expected after an R\&R.  Items 6--8 are polish.

\vspace{0.5em}
\hrule
\vspace{0.5em}

\section{5.\ JBES-Specific Considerations}

\subsection{5.1\quad Scope and Fit}

JBES publishes methodological contributions oriented toward applications in
micro- and macroeconomics.  It regularly publishes Bayesian methods papers.
The paper's balance of methodology (variance decomposition, cumulant matching)
and application (RDU estimation with real experimental data) is an excellent fit.

JBES values:
\begin{itemize}[leftmargin=1.5em]
\item Practical applicability of methods to real economic problems
\item Clear notation and exposition (the paper reads well)
\item Demonstrable consequences for inference (the CE analysis does this)
\item Engagement with existing econometric methodology
\end{itemize}

\subsection{5.2\quad Expected Length}

JBES has no strict word limit but typical articles are 15--25 pages.  The
current paper is roughly 14--16 pages of body text (including appendices).
With the additions above (literature review $+$ sensitivity analysis $+$ tables),
the paper would reach 20--24 pages---comfortably within JBES norms.

\subsection{5.3\quad Formatting Requirements}

\begin{itemize}[leftmargin=1.5em]
\item JBES uses Chicago-style author-year citations.  The current
\texttt{IEEEtran} + \texttt{authoryear-comp} biblatex style may need
adjustment.
\item JBES requires a data/code availability statement.
\item Double-spacing is already in place.
\item Supplementary materials (appendices, additional tables/figures) can be
submitted as online appendices.
\end{itemize}

\subsection{5.4\quad Positioning for JBES Audience}

For JBES, the framing should emphasize:
\begin{itemize}[leftmargin=1.5em]
\item This is a \emph{practical tool} for applied researchers, not a theoretical
exercise.  The variance decomposition reduces $4p$ difficult hyperprior choices
to 2 interpretable ones.
\item The method applies beyond risk preferences---any BHM with
Normal-Inverse-Gamma (or NIW) hyperpriors.
\item The comparison with Gao et al.\ (2022) shows that hyperprior choice has
substantive consequences---not just for convergence, but for economic inference
(the CE shifts are economically meaningful).
\item The reparameterization has the attractive property of being unit-agnostic,
which is important for multidimensional models where parameters live on
different scales.
\end{itemize}

\vspace{0.5em}
\hrule
\vspace{0.5em}

\section{6.\ What Has Been Accomplished}

For reference, here is a summary of the progress against the original 8-item
action list:

\begin{center}
\small
\begin{tabular}{@{}clcc@{}}
\toprule
\textbf{\#} & \textbf{Item} & \textbf{Status} & \textbf{Priority} \\
\midrule
1 & Literature review and positioning & Not started & Critical \\
2 & Sensitivity analysis & Not started & Critical \\
3 & Complete Appendix~A & \textbf{Done} & --- \\
4 & Posterior summary table & Not started & High \\
5 & Expand uninformative comparison & Not started & High \\
6 & Simulation study & Not started & Medium \\
7 & Figure captions and typos & Partial & Low \\
8 & Discussion restructuring & Not started & Medium \\
\midrule
  & Expand Appendix~C (MCMC details) & \textbf{Done} & --- \\
  & Prose consistency (we $\to$ I) & \textbf{Done} & --- \\
\bottomrule
\end{tabular}
\end{center}

\vspace{0.5em}
\hrule
\vspace{0.5em}

\section{7.\ Suggested Reading List}

Unchanged from the initial review.  Papers marked with ${}^{\star}$ are
essential reading before revising.

\subsection{7.1\quad Must-Read (${}^{\star}$)}

\begin{enumerate}[leftmargin=1.5em]
\item[${}^{\star}$] Gelman, A.\ (2006).  Prior distributions for variance
parameters in hierarchical models.  \textit{Bayesian Analysis}, 1(3), 515--534.

\item[${}^{\star}$] Simpson, D., Rue, H., Riebler, A., Martins, T.\,G., \&
S{\o}rbye, S.\,H.\ (2017).  Penalising model component complexity.
\textit{Statistical Science}, 32(1), 1--28.

\item[${}^{\star}$] Polson, N.\,G.\ \& Scott, J.\,G.\ (2012).  On the
half-Cauchy prior for a global scale parameter.  \textit{Bayesian Analysis},
7(4), 887--902.

\item[${}^{\star}$] Gabry, J., Simpson, D., Vehtari, A., Betancourt, M., \&
Gelman, A.\ (2019).  Visualization in Bayesian workflow.  \textit{JRSS-A},
182(2), 389--402.

\item[${}^{\star}$] Roos, M., Martins, T.\,G., Held, L., \& Rue, H.\ (2015).
Sensitivity analysis for Bayesian hierarchical models.  \textit{Bayesian
Analysis}, 10(2), 321--349.
\end{enumerate}

\subsection{7.2\quad Strongly Recommended}

\begin{enumerate}[leftmargin=1.5em]
\item[6.] Berger, J.\,O., Strawderman, W., \& Tang, D.\ (2005).  Posterior
propriety and admissibility of hyperpriors in normal hierarchical models.
\textit{Annals of Statistics}, 33(2), 606--646.

\item[7.] Huang, A.\ \& Wand, M.\,P.\ (2013).  Simple marginally
noninformative prior distributions for covariance matrices.  \textit{Bayesian
Analysis}, 8(2), 439--452.

\item[8.] Natarajan, R.\ \& Kass, R.\,E.\ (2000).  Reference Bayesian methods
for generalized linear mixed models.  \textit{JASA}, 95(449), 227--237.

\item[9.] Daniels, M.\,J.\ \& Kass, R.\,E.\ (1999).  Nonconjugate Bayesian
estimation of covariance matrices.  \textit{JASA}, 94, 1254--1263.

\item[10.] Nilsson, H., Rieskamp, J., \& Wagenmakers, E.-J.\ (2011).
Hierarchical Bayesian parameter estimation for cumulative prospect theory.
\textit{J.\ Mathematical Psychology}, 55, 84--93.

\item[11.] Mikkola, P.\ et al.\ (2023).  Prior knowledge elicitation: The past,
present, and future.  \textit{Bayesian Analysis}, advance publication.

\item[12.] Vehtari, A., Gelman, A., \& Gabry, J.\ (2017).  Practical Bayesian
model evaluation using leave-one-out cross-validation and WAIC.
\textit{Statistics and Computing}, 27, 1413--1432.
\end{enumerate}

\subsection{7.3\quad Additional Context}

\begin{enumerate}[leftmargin=1.5em]
\item[13.] Conte, A., Hey, J.\,D., \& Moffatt, P.\,G.\ (2011).  Mixture models
of choice under risk.  \textit{J.\ Econometrics}, 162(1), 79--88.

\item[14.] Hey, J.\,D.\ \& Orme, C.\ (1994).  Investigating generalizations of
expected utility theory using experimental data.  \textit{Econometrica}, 62(6),
1291--1326.

\item[15.] Betancourt, M.\ \& Girolami, M.\ (2015).  Hamiltonian Monte Carlo for
hierarchical models.  In \textit{Current Trends in Bayesian Methodology with
Applications}, 79--102.

\item[16.] Browne, W.\,J.\ \& Draper, D.\ (2006).  A comparison of Bayesian and
likelihood-based methods for fitting multilevel models.  \textit{Bayesian
Analysis}, 1(3), 473--514.

\item[17.] Chung, Y., Rabe-Hesketh, S., Dorie, V., Gelman, A., \& Liu, J.\
(2013).  A nondegenerate penalized likelihood estimator for variance parameters
in multilevel models.  \textit{Psychometrika}, 78(4), 685--709.

\item[18.] Gelman, A.\ \& Hill, J.\ (2007).  \textit{Data Analysis Using
Regression and Multilevel/Hierarchical Models}.  Cambridge University Press.

\item[19.] Poirier, D.\,J.\ (1995).  \textit{Intermediate Statistics and
Econometrics: A Comparative Approach}.  MIT Press.  (Chapter~10 on prior
sensitivity.)

\item[20.] Kass, R.\,E.\ \& Raftery, A.\,E.\ (1995).  Bayes factors.
\textit{JASA}, 90(430), 773--795.
\end{enumerate}

\vspace{0.5em}
\hrule
\vspace{0.5em}

\section{8.\ Summary Assessment}

\textbf{Previous assessment:} Reject (incomplete appendix, insufficient
literature engagement, no sensitivity analysis).

\textbf{Current assessment:} The technical core is now complete and solid.  The
paper is not yet ready for JBES submission, but the remaining gaps are
addressable.

\textbf{Path to JBES submission:}
\begin{itemize}[leftmargin=1.5em]
\item \textbf{Minimum viable submission} (items 1--3): Literature positioning,
sensitivity analysis, posterior summary table.  This brings the paper to
${\sim}$20 pages and addresses the most likely desk-reject triggers.

\item \textbf{Competitive submission} (items 1--5): Add a ``reasonable'' prior
comparator and either a simulation study or second application.  This brings
the paper to ${\sim}$22--24 pages and gives referees the robustness evidence
they will want.

\item \textbf{Strong submission} (all items): Full polish including discussion
restructuring, prior specification guidance, figure captions, and bibliography
alignment.  This is a paper that should receive an R\&R.
\end{itemize}

\textbf{If JBES declines:} The paper can be submitted to \textit{Econometric
Reviews} (Brief Report format, 15 pages) or \textit{Bayesian Analysis} with
minimal changes.  The completed appendices give it the technical depth for
either venue.

The core idea---specify a prior for an individual, derive hyperpriors via
variance decomposition so the unconditional BHM prior matches---is clean,
practical, and potentially impactful.  With the literature positioning and
sensitivity analysis, this is a competitive JBES submission.

\end{document}
